 %%%%%%%%%%%%%%%%%%%%%%%%%%%%%%%%%%%%%%%%%
% Beamer Presentation
% LaTeX Template
% Version 1.0 (10/11/12)
%
% This template has been downloaded from:
% http://www.LaTeXTemplates.com
%
% License:
% CC BY-NC-SA 3.0 (http://creativecommons.org/licenses/by-nc-sa/3.0/)
%
%%%%%%%%%%%%%%%%%%%%%%%%%%%%%%%%%%%%%%%%%

%----------------------------------------------------------------------------------------
%	PACKAGES AND THEMES
%----------------------------------------------------------------------------------------

\documentclass[9pt,aspectratio=169]{beamer}
\usepackage[utf8]{inputenc}
\usepackage{bm}
\usepackage{bbm}
\usepackage{graphicx}
\usepackage{cmbright}
\usefonttheme{professionalfonts}
\DeclareSymbolFont{greekletters}{OML}{sansmath}{m}{it}
\DeclareMathSymbol{\beta}{\mathord}{greekletters}{"0C}
%\usefonttheme{serif}
% The Beamer class comes with a number of default slide themes
% which change the colors and layouts of slides. Below this is a list
% of all the themes, uncomment each in turn to see what they look like.

%\usetheme{default}
%\usetheme{AnnArbor}
% \usetheme{Antibes}
%\usetheme{Bergen}
%\usetheme{Berkeley}
%\usetheme{Berlin}
%\usetheme{Boadilla}
% \usetheme{CambridgeUS}
%\usetheme{Copenhagen}
%\usetheme{Darmstadt}
% \usetheme{Dresden}
\usetheme{Frankfurt}
%\usetheme{Goettingen}
% \usetheme{Hannover}
% \usetheme{Ilmenau}
% \usetheme{JuanLesPins}
%\usetheme{Luebeck}
%\usetheme{Madrid}
%\usetheme{Malmoe}
%\usetheme{Marburg}
% \usetheme{Montpellier}
% \usetheme[width=2.5cm]{PaloAlto}
%\usetheme{Pittsburgh}
% \usetheme{Rochester}
%\usetheme{Singapore}
% \usetheme{Szeged}
% \usetheme{Warsaw}
% \setbeamertemplate{title page}[default][colsep=-4bp,rounded=false]
\setbeamertemplate{headline}{}
\defbeamertemplate{footline}{myframe number}
{
  \hspace{1pt}
    \usebeamercolor[structure]{page number in head/foot}%
  \usebeamerfont{page number in head/foot}%
  \raisebox{1.7ex}[0pt][0pt]{% <--- change here
    \insertframenumber\,/\,\inserttotalframenumber\kern1em}%
}
\setbeamertemplate{footline}[myframe number]
\AtBeginSection[]
{
  \begin{frame}
    \frametitle{Outline}
    \begin{minipage}{1.0\linewidth}
      \tableofcontents[currentsection,currentsubsection]
    \end{minipage}
  \end{frame}
}

\AtBeginSubsection[]
{
  \begin{frame}
    \frametitle{Outline}
    \begin{minipage}{1.0\linewidth}
      \tableofcontents[currentsection,currentsubsection]
    \end{minipage}
  \end{frame}
}

% As well as themes, the Beamer class has a number of color themes
% for any slide theme. Uncomment each of these in turn to see how it
% changes the colors of your current slide theme.
%\definecolor{mycolor}{rgb}{.8,.0,.3}
%\setbeamercolor*{palette primary}{use=structure,fg=white,bg=mycolor}
%\usecolortheme{albatross}
%\usecolortheme{beaver}
%\usecolortheme{beetle}
%\usecolortheme{crane}
%\usecolortheme{dolphin}
%\usecolortheme{dove}
%\usecolortheme{fly}
%\usecolortheme{lily}
%\usecolortheme{orchid}
% \usecolortheme{rose}
%\usecolortheme{seagull}
\usecolortheme{seahorse}
%\usecolortheme{whale}
%\usecolortheme{wolverine}
%\usecolortheme[named=mycolor]{structure}
%\setbeamertemplate{footline} % To remove the footer line in all slides uncomment this line
%\setbeamertemplate{footline}[page number] % To replace the footer line in all slides with a simple slide count uncomment this line

%\setbeamertemplate{navigation symbols}{} % To remove the navigation symbols from the bottom of all slides uncomment this line


\usepackage{booktabs} % Allows the use of \toprule, \midrule and \bottomrule in tables

% Custom commands
\newcommand{\R}{\mathbb{R}}
\newcommand{\Z}{\mathbb{Z}}

\title{Lecture 7: From Static to Dynamic Optimization}
\date{}

\author{Junnan Zhang} % Your name
\institute[Paula and Gregory Chow Institute for Studies in Economics\\
Xiamen University] % Your institution as it will appear on the bottom of every slide, may be shorthand to save space
{
	Paula and Gregory Chow Institute for Studies in Economics \\ Xiamen University  \\ % Your institution for the title page
	\medskip
}
\date{Fall, 2026} % Date, can be changed to a custom date
\setbeamertemplate{itemize items}[default]
\setbeamertemplate{sidebar}[sidebar right]
% \definecolor{myco}{HTML}{a67678}
\definecolor{myco}{HTML}{8776a6}
\setbeamercolor{itemize item}{fg=myco} % all frames will have red bullets
\setbeamercolor{itemize subitem}{fg=myco} % all frames will have red bullets
\setbeamercolor{block title}{bg=myco}
\setbeamercolor{structure}{fg=myco}
% \setbeamercovered{transparent}


\begin{document}

\frame{\titlepage}


\begin{frame}{Overview}
    \begin{itemize}
        \item Brief introduction to finite-horizon dynamic optimization in
        discrete time under certainty
        \item Introduce the Bellman equation and the Principle of Optimality
    \end{itemize}
\end{frame}

\section{Problem Description}

\begin{frame}
    \frametitle{Problem Description}
    A finite-horizon dynamic optimization problem
    \begin{align*}
        \sup_{\{x_t\}_{t=1}^{T}} &\sum_{t=0}^{T-1}F_t(x_t, x_{t+1})\\
        \text{s.t. }
        &x_{t+1} \in \Gamma_t(x_t) \text{ for all } t = 0, 1, \ldots, T-1\\
        &x_0 \text{ given}
    \end{align*}
    \small
    \begin{itemize}
        \item $x_t \in X$ is the state variable (state vector)
        \item $F_t:X \times X \to \R$ is the instantaneous payoff function
        \item $\Gamma_t(x)$ is a non-empty set-valued mapping (correspondence):
        $\Gamma_t: X \rightrightarrows X$
        \item The constraint $x_{t+1} \in \Gamma_t(x_t)$ captures
        feasibility: tomorrow's state must be reachable from today's state
        \item We use $\sup$ instead of $\max$ because the maximum may not be
        attained in general settings
    \end{itemize}
\end{frame}

\begin{frame}
    \frametitle{Example: Cake Eating Problem}
    \begin{itemize}
        \item Finite time horizon: $t = 0, 1, \ldots, T$
        \item At $t=0$ the agent is given a complete cake with size $\bar x$
        \item Let $x_t$ denote the size of the cake at the beginning of each
        period, so that, in particular, $x_0=\bar x$
        \item After choosing to consume $c_t$ of the cake in period $t$ there is
        $$
        x_{t+1} = x_t - c_t
        $$
        left in period $t+1$
        \item Consuming quantity $c$ of the cake gives current utility
        $u(c)$
    \end{itemize}
\end{frame}

\begin{frame}
    \frametitle{Example: Cake Eating Problem}
    \begin{itemize}
        \item The agent's problem can be written as
        \begin{equation*}
            \max_{\{c_t\}} \sum_{t=0}^{T-1} \beta^t u(c_t)            
        \end{equation*}
        subject to
        \begin{equation*}
            x_{t+1} = x_t - c_t \quad 0\leq c_t\leq x_t, \quad x_0 = \bar{x}
        \end{equation*}
        \item Formulate this in the general framework:
        \begin{itemize}
            \item $F_t(x_t, x_{t+1}) = \beta^t u(x_t - x_{t+1})$
            \item $\Gamma_t(x_t) = [0, x_t]$
            \item We write $c_t = x_t - x_{t+1}$ in terms of the states
            \item The feasible set $\Gamma_t(x_t) = [0, x_t]$ encodes both
            non-negativity and resource constraints
        \end{itemize}
    \end{itemize}
\end{frame}

\begin{frame}
    \frametitle{Example: Cake Eating Problem}
    The key trade-off in the cake-eating problem is:
    \begin{itemize}
        \item Delaying consumption is costly because of the discount factor
        \item But delaying some consumption is also attractive because $u$ is
        concave. The concavity of $u$ implies that the consumer gains value
        from consumption smoothing, which means spreading consumption out
        over time
        \item The optimal solution balances these opposing forces
    \end{itemize}
\end{frame}

\begin{frame}
    \frametitle{Lagrangian Approach?}
    \begin{itemize}
        \item In this example, the constraint $x_{t+1} \in \Gamma_t(x_t)$ can
        be written as $x_{t+1} \geq 0$ and $x_{t+1} \leq x_t$
        \item The Lagrangian $L(x_1, x_2, \ldots, x_T, \lambda_1, \ldots,
        \lambda_T, \mu_1, \ldots, \mu_T)$
        \item We need to solve for $3T$ unknowns from the first order
        conditions: $T$ state variables and $2T$ multipliers
        \item As $T$ grows, this becomes a large simultaneous system whose
        dimension grows with $T$
        \item Dynamic programming offers a more efficient approach by
        breaking the problem into stages
    \end{itemize}
\end{frame}

\begin{frame}
    \frametitle{Lagrangian Approach?}
    \begin{itemize}
        \item Consider an example where $T = 3$
        \item Adopt the CRRA utility function:
        \begin{equation*}
            u(c) = \frac{c^{1-\gamma}}{1-\gamma} \qquad (0 < \gamma < 1)
        \end{equation*}
        \item CRRA = Constant Relative Risk Aversion; $\gamma$ measures
        the degree of risk aversion
        \item Higher $\gamma$ means stronger preference for consumption smoothing
        \item The case $\gamma = 1$ corresponds to $u(c) = \ln(c)$
        (logarithmic utility)
        \item Formulate the Lagrangian
    \end{itemize}
\end{frame}

\section{Recursive Formulation and Dynamic Programming}

\begin{frame}
    \frametitle{Dynamic Programming}
    \begin{itemize}
        \item Given the recursive structure of this problem, we introduce
        the \textbf{value function} $V_t(x)$:
        \begin{align*}
            V_t(x) &:= \sup_{\{x_s\}_{s=t+1}^{T}} \sum_{s=t}^{T-1}F_s(x_s, x_{s+1})\\
            \text{s.t. } &x_{s+1} \in \Gamma_s(x_s) \text{ for all } s = t,
            \ldots, T-1\\
            &x_t = x
        \end{align*}
        \item $V_t(x)$ gives the lifetime value at time $t$ if the state is
        $x_t = x$: ``what is the best I can do from time $t$ onward, starting
        with state $x$?''
        \item This is the \emph{value of the problem} viewed from period $t$
        with initial condition $x_t = x$
        \item Note that $V_0(x_0)$ is the value of the original problem
    \end{itemize}
\end{frame}

\begin{frame}
    \frametitle{Dynamic Programming Principle}
    \begin{block}{Dynamic Programming Principle}
        If $(x_1, x_2, \ldots, x_T)$ is a solution to problem $V_0(x_0)$,
        then for any $t = 1, 2, \ldots, T-1$, $(x_{t+1}, \ldots, x_T)$ is a
        solution to problem $V_t(x_t)$.
    \end{block}
    Intuition: an optimal path remains optimal from any point along it.
    
    Proof by contradiction:
    \begin{itemize}
        \item Suppose that there is an admissible sequence $(z_{t+1}, \ldots,
        z_T)$ that achieves higher value than $(x_{t+1}, \ldots, x_T)$ at
        time $t$ when the state is $x_t$
        \item Show that $(x_0, x_1, \ldots, x_t, z_{t+1}, \ldots, z_T)$ is
        admissible and achieves higher value of $V_0(x_0)$
        \item In other words, we could "splice" this better tail onto our
        path, improving overall value, contradicting optimality
    \end{itemize}
\end{frame}

\begin{frame}
    \frametitle{Bellman Equation}
    \begin{block}{Bellman Equation}
        For all $t = 0, 1, \ldots, T-1$, we have
        \begin{equation*}
            V_t(x) = \sup_{x' \in \Gamma_t(x)} \left\{F_t(x, x') +
                V_{t+1}(x')\right\},
        \end{equation*}
        where we use the terminal convention $V_T(x)=0$.
    \end{block}
    Interpretation: today's value = today's payoff + future value.
    \begin{itemize}
        \item Choose the next period state $x'$ to maximize: current return
        $F_t(x, x')$ plus continuation value $V_{t+1}(x')$
        \item This reduces a $T$-period problem to a sequence of 2-period
        problems
        \item $V_{t+1}(x')$ summarizes all future consequences of choosing
        $x'$ today
        \item Note that we do not assume the existence of an optimal sequence here
    \end{itemize}
\end{frame}

\begin{frame}
    \frametitle{Bellman Equation: Proof Sketch}
    Fix $t=0,1,\ldots,T-1$ and a state $x$.
    \begin{itemize}
        \item First, we prove that $V_t(x) \geq \sup_{x' \in \Gamma_t(x)}
        \left\{ F_t(x, x') + V_{t+1}(x')\right\}$
        \begin{enumerate}
            \item For any admissible sequence $(x, x_{t+1}, \ldots, x_T)$,
            show that
            \begin{equation*}
                V_t(x) \geq F_t(x, x_{t+1}) + \sum_{s=t+1}^{T-1} F_s(x_s, x_{s+1})
            \end{equation*}
            \item Taking the supremum over $(x_{t+2}, \ldots, x_T)$ gives
            $V_t(x) \geq F_t(x, x_{t+1}) + V_{t+1}(x_{t+1})$
            \item Taking the supremum again over $x_{t+1}$ proves the result
        \end{enumerate}
        \item Next, we prove that $V_t(x) \leq \sup_{x' \in \Gamma_t(x)}
        \left\{ F_t(x, x') + V_{t+1}(x')\right\}$
        \begin{enumerate}
            \item For any $\epsilon > 0$, let $(x, x_{t+1}, \ldots, x_T)$
            be an admissible sequence such that
            \begin{equation*}
                V_t(x) \leq F_t(x, x_{t+1}) + \sum_{s=t+1}^{T-1} F_s(x_s,
                x_{s+1}) + \epsilon
            \end{equation*}
            \item Show that $\sum_{s=t+1}^{T-1} F_s(x_s, x_{s+1}) \leq
            V_{t+1}(x_{t+1})$
            \item Taking the supremum over $x_{t+1}$ and taking the limit as
            $\epsilon \to 0$ completes the proof
        \end{enumerate}
    \end{itemize}
\end{frame}

\begin{frame}
    \frametitle{Backward Induction}
    \begin{itemize}
        \item Using the Bellman equation, we can compute the value function
        through backward induction.
        \item In the last period, $V_{T-1}(x) = \sup_{x' \in \Gamma_{T-1}(x)}
        F_{T-1}(x, x')$. This is a static optimization problem. Solving this
        problem gives: (1) the value at $T-1$ as a function of $x_{T-1}$, and
        (2) the optimal $x_T$ as a function of $x_{T-1}$.
        \item Then, for each $t = 0, 1, \ldots, T-2$, we can solve for
        $V_t(x)$ and obtain the optimal $x_{t+1}$ as a function of $x_t$:
        \begin{equation*}
            V_{t}(x) = \sup_{x' \in \Gamma_t(x)} \left\{ F_t(x, x')
                + V_{t+1}(x') \right\}
        \end{equation*}
        \item Given $x_0$, we can also obtain an admissible sequence by
        solving these Bellman equations.
        \item The question is: is this sequence our solution?
    \end{itemize}
\end{frame}

\begin{frame}
    \frametitle{Principle of Optimality}
    \begin{block}{Principle of Optimality}
        The sequence $(x_1, \ldots, x_T)$ is a solution to $V_0(x_0)$ if
        and only if for all $t = 0, 1, \ldots, T-1$, $x_{t+1}$ is a solution
        to
        \begin{equation*}
            \sup_{x' \in \Gamma_t(x_t)} \left\{ F_t(x_t, x') + V_{t+1}(x') \right\},
        \end{equation*}
        where $V_T=0$.
    \end{block}
    This is the characterization theorem: it tells us exactly when a sequence
    is optimal

\end{frame}

\begin{frame}
    \frametitle{Principle of Optimality: Proof Sketch}
    Proof:
    \begin{itemize}
        \item Let $(x_1, \ldots, x_T)$ be a solution to $V_0(x_0)$. The
        dynamic programming principle implies that $V_t(x_t) =
        \sum_{s=t}^{T-1} F_s(x_s, x_{s+1})$ for all $t=0,1,\ldots,T-1$.
        By definition of the value function and Bellman equation:
        \begin{equation*}
            F_t(x_t, x_{t+1}) + V_{t+1}(x_{t+1}) = V_{t}(x_t) = \sup_{x' \in
              \Gamma_t(x_t)} \left\{ F_t(x_t, x') + V_{t+1}(x') \right\}
        \end{equation*}
        \item Let $x_{t+1}$ be a solution to $\sup_{x' \in \Gamma_t(x_t)}
        \left\{ F_t(x_t, x') + V_{t+1}(x') \right\}$ for all $t = 0, 1,
        \ldots, T-1$. Using $V_T(x_T)=0$, iterating on the Bellman equation gives
        \begin{equation*}
            V_0(x_0) = F_0(x_0, x_1) + V_1(x_1) = F_0(x_0, x_1) + F_1(x_1,
            x_2) + V_2(x_2) = \ldots = \sum_{t=0}^{T-1} F_t(x_t, x_{t+1})
        \end{equation*}
    \end{itemize}
\end{frame}

\section{Alternative Formulation}

\begin{frame}
    \frametitle{Problem Description}
    An alternative (more general) formulation
    \begin{align*}
        \sup_{\{u_t\}_{t=0}^{T-1}} &\sum_{t=0}^{T-1}F_t(x_t, u_t)\\
        \text{s.t. }
        &u_t \in G_t(x_t) \\
        &x_{t+1} = f_t(x_t, u_t)\\
        &x_0 \text{ given}
    \end{align*}
    \small
    \begin{itemize}
        \item $u_t \in U$ is the control variable
        \item $F_t:X \times U \to \R$ is the instantaneous payoff function
        \item $G_t(x)$ is the feasible correspondence:
        $G_t: X \rightrightarrows U$
        \item $f_t: X \times U \to X$ gives the next period state
    \end{itemize}
\end{frame}

\begin{frame}
    \frametitle{Value Function}
    The value function is
    \begin{align*}
        V_t(x) &:= \sup_{\{u_s\}_{s=t}^{T-1}} \sum_{s=t}^{T-1} F_s(x_s,u_s),\\
        \text{s.t. } &u_s \in G_s(x_s) \text{ for all } s = t,
        \ldots, T-1\\
        &x_{s+1}=f_s(x_s,u_s)\text{ for all }s=t,\ldots,T-1\\
        &x_t = x
    \end{align*}
\end{frame}

\begin{frame}
    \frametitle{Dynamic Programming Theorems}
    \begin{block}{Dynamic Programming Principle}
        If a control sequence $(u_0, u_1, \ldots, u_{T-1})$ with induced
        states $(x_1, x_2, \ldots, x_T)$ is a solution to problem $V_0(x_0)$,
        then for any $t = 1, 2, \ldots, T-1$, $(u_t, \ldots, u_{T-1})$ is a
        solution to problem $V_t(x_t)$.
    \end{block}
    
    \begin{block}{Bellman Equation}
        For all $t = 0, 1, \ldots, T-1$, we have
        \begin{equation*}
            V_t(x) = \sup_{u \in G_t(x)} \left\{F_t(x, u) + V_{t+1}(f_t(x,
                u))\right\}
        \end{equation*}
    \end{block}

    \begin{block}{Principle of Optimality}
        The control sequence $(u_0, \ldots, u_{T-1})$ with induced states
        $(x_1, x_2, \ldots, x_T)$ is a solution to $V_0(x_0)$ if and only if
        for all $t = 0, 1, \ldots, T-1$, $u_t$ is a solution to
        \begin{equation*}
            \sup_{u \in G_t(x_t)} \left\{ F_t(x_t, u) + V_{t+1}(f_t(x_t, u))
            \right\}
        \end{equation*}
        where $x_{t+1} = f_t(x_t, u_t)$.
    \end{block}
\end{frame}


\section{Examples}

\begin{frame}[allowframebreaks]
    \frametitle{Cake Eating Problem}
    \begin{itemize}
        \item Consider again the cake eating problem with CRRA utility and $T
        = 3$
        \item The Bellman equation is
        \begin{equation*}
            V_t(x_t) = \max_{x_{t+1} \in [0, x_t]} \left\{ \beta^t\frac{(x_t -
                  x_{t+1})^{1-\gamma}}{1-\gamma} + V_{t+1}(x_{t+1})
            \right\}
        \end{equation*}
        \item We use $\max$ here because the optimum exists (continuous
        function on compact set)
        \item Solve it via backward induction
    \end{itemize}
    \framebreak
    \begin{itemize}
        \item When $t = 2$:
        \begin{equation*}
            V_2(x_2) = \max_{x_3 \in [0, x_2]} \beta^2\frac{(x_2 -
              x_3)^{1-\gamma}}{1-\gamma} = \beta^2 \frac{x_2^{1-\gamma}}{1-\gamma}
        \end{equation*}
        and the optimal policy is $x_3 = 0$: eat everything remaining (no
        reason to save)
        \item When $t = 1$:
        \begin{equation*}
            V_1(x_1) = \max_{x_2 \in [0, x_1]}\left\{\beta \frac{(x_1 -
                  x_2)^{1-\gamma}}{1-\gamma} + \beta^2
                \frac{x_2^{1-\gamma}}{1-\gamma} \right\} =
            \frac{x_1^{1-\gamma}}{1-\gamma} \frac{\beta}{(1 +
              \beta^{\frac{1}{\gamma}})^{-\gamma}}
        \end{equation*}
        with $x_2 = x_1/(1 + \beta^{-\frac{1}{\gamma}})$
        \item When $t = 0$:
        \begin{equation*}
            V_0(x_0) = \max_{x_1 \in [0, x_0]}\left\{ \frac{(x_0 -
                  x_1)^{1-\gamma}}{1-\gamma} + \frac{\beta}{(1 +
                  \beta^{\frac{1}{\gamma}})^{-\gamma}}
                \frac{x_1^{1-\gamma}}{1-\gamma} \right\}
        \end{equation*}
        which can be solved although a bit tedious
    \end{itemize}
\end{frame}


\begin{frame}{Exercises}
    \begin{itemize}
        \item Let $\Gamma_0(x_0) := [0, x_0^4 + 2x_0 + 3]$, $\Gamma_1(x_1) :=
        \left[\frac{x_1}{2}, x_1^2 + x_1\right]$, and $\Gamma_2(x_2) :=
        \left[0, \frac{x_2^2 + 4}{x_2^2}\right]$
        \item Define $f(x_1, x_0) := 2x_1x_0 - x_1^2 + x_1$, $g(x_2, x_1) :=
        -\frac{1}{2x_2} + x_2x_1 - \frac{x_2^2}{2}$, $h(x_3, x_2) :=
        \sqrt{x_3} - \frac{x_3x_2}{2}$
        \item Consider the following problem
        \begin{align*}
            \sup_{x_1,x_2,x_3} &\left\{f(x_1, x_0) + g(x_2, x_1) + h(x_3, x_2)
            \right\}\\
            \text{s.t. } &x_i \in \Gamma_{i-1}(x_{i-1}) \text{ for all } i =
            1, 2, 3\\
            & x_0 \geq 0 \text{ given}
        \end{align*}
        \item This is a 3-stage problem with state-dependent constraint sets
        \item The payoff functions and feasibility constraints have different
        functional forms at each stage
    \end{itemize}
\end{frame}

\begin{frame}{Exercise}
    Stage 3:
    \begin{itemize}
        \item Compute $V_2(x_2)$
        \[
            V_2(x_2) = \sup_{y \in \Gamma_2(x_2)} h(y, x_2) = \sup_{y \in
              \Gamma_2(x_2)} \left\{\sqrt{y} - \frac{yx_2}{2}\right\}
        \]
        \item First-order condition:
        \[
            \frac{d}{d y}\left(\sqrt{y} - \frac{yx_2}{2}\right)
            = \frac{1}{2\sqrt{y}} - \frac{x_2}{2} = 0
        \]
        \item Solving:
        \[
            \frac{1}{2\sqrt{y}} = \frac{x_2}{2} \implies y^* = \frac{1}{x_2^2}
        \]
        \item Check second order condition: $-\frac{1}{4y^{3/2}} < 0$, so
        this is indeed a maximum
        \item Since $y^* \in \Gamma_2(x_2)$, $V_2(x_2) = {\frac{1}{2x_2}}$
    \end{itemize}
\end{frame}

\begin{frame}{Exercise}
    Stage 2:
    \begin{itemize}
        \item Compute $V_1(x_1)$:
        \begin{align*}
            V_1(x_1) &= \sup_{y \in \Gamma_1(x_1)} \left\{g(y, x_1) + V_2(y)\right\}\\
            &= \sup_{y \in \Gamma_1(x_1)} \left\{-\frac{1}{2y} + yx_1 -
                \frac{y^2}{2} + \frac{1}{2y}\right\}
        \end{align*}
        \item First-order condition gives $y^* = x_1$
        \item Second-order condition is $-1 < 0$
        \item Check feasibility: $y^* = x_1$ must satisfy $\frac{x_1}{2} \leq
        x_1 \leq x_1^2 + x_1$, which holds for $x_1 \geq 0$
        \item The value function is $V_1(x_1) = x_1^2/2$
    \end{itemize}
\end{frame}

\begin{frame}{Exercise}
    Stage 1:
    \begin{itemize}
        \item Compute $V_0(x_0)$:
        \begin{equation*}
            V_0(x_0) = \sup_{y \in \Gamma_0(x_0)} \left\{2yx_0 + y -
                \frac{y^2}{2}\right\}
        \end{equation*}
        \item First-order condition gives $y = 2x_0 + 1$
        \item Second-order condition: $-1 < 0$
        \item Check feasibility: $2x_0 + 1 \in [0, x_0^4 + 2x_0 + 3]$
        \item The value function is $V_0(x_0) = 2x_0^2 + 2x_0 + 1/2$
    \end{itemize}
    The optimal policy is:
    \begin{equation*}
        \left(2x_0 + 1, 2x_0 + 1, \frac{1}{(2x_0 + 1)^2}\right)        
    \end{equation*}
\end{frame}

\end{document}
